% !TEX root = ../../main.tex
% Shared physical dimensions for interval and number-line diagrams.
\tikzset{
  interval diagram/.style={
    x=1cm,y=6mm,font=\footnotesize,
    color=black,line width=0.5pt,line cap=round,line join=round,
    >={Stealth[length=4pt,width=3pt]}
  },
  interval endpoint/.style={
    circle,draw=black,line width=0.5pt,
    minimum size=4pt,inner sep=0pt,outer sep=0pt
  },
  interval open/.style={interval endpoint,fill=white},
  interval closed/.style={interval endpoint,fill=black},
  interval label/.style={inner sep=1pt}
}

\begin{tikzpicture}[interval diagram,x=8mm,y=10mm]
  % Successive halves retain infinitely many indexed sequence terms.
  \draw[gray,dashed] (1.72,0.5) -- (1.72,-3.3);
  \node[above=5pt] at (1.72,0.5) {$L$};

  \draw (0,0) -- (4,0);
  \node[interval closed] at (0,0) {};
  \node[interval closed] at (4,0) {};
  \node[left=6pt] at (0,0) {$I_1$};
  \fill (3.5,0) circle[radius=1.4pt];
  \node[above=5pt] at (3.5,0) {$a_{n_1}$};

  \draw (0,-1.4) -- (2,-1.4);
  \node[interval closed] at (0,-1.4) {};
  \node[interval closed] at (2,-1.4) {};
  \node[left=6pt] at (0,-1.4) {$I_2$};
  \fill (0.45,-1.4) circle[radius=1.4pt];
  \node[above=5pt] at (0.45,-1.4) {$a_{n_2}$};

  \draw (1,-2.8) -- (2,-2.8);
  \node[interval closed] at (1,-2.8) {};
  \node[interval closed] at (2,-2.8) {};
  \node[left=6pt] at (0,-2.8) {$I_3$};
  \fill (1.28,-2.8) circle[radius=1.4pt];
  \node[above left=5pt] at (1.28,-2.8) {$a_{n_3}$};
  \node at (1.4,-3.8) {$\vdots$};
\end{tikzpicture}
